\documentclass[10pt]{article}
\usepackage[margin=0.83in]{geometry}
\usepackage{amsmath,amssymb,booktabs}
\usepackage[T1]{fontenc}
\usepackage{lmodern}
\usepackage{microtype}
\usepackage[colorlinks=true,linkcolor=black,urlcolor=blue]{hyperref}
\setlength{\parindent}{0pt}
\setlength{\parskip}{5pt plus 1pt}
\setlength{\abovedisplayskip}{7pt plus 2pt}
\setlength{\belowdisplayskip}{7pt plus 2pt}
\renewcommand{\arraystretch}{1.12}

\title{Closing the Count Gap in Quicksort\\
\large What a smooth law can match, and what the tree still knows}
\author{}
\date{October 2, 2026}

\begin{document}
\maketitle
\vspace{-1.8em}

\begin{abstract}
Randomized Quicksort has a discrete comparison-count law at every finite size,
but its limiting law is continuous. We formalize a way to measure the finite-size
gap in bits, give a construction that closes it exactly, and explain why that
construction does not recover the tree's dynamics. A finite-size split model and
an error bound show how to pursue a useful, compact approximation instead.
\end{abstract}

\section*{From a good curve to an evolving model}
Let $C_n$ count comparisons when uniform-pivot Quicksort sorts $n$ distinct
keys. Its exact law obeys
\begin{equation}
 C_n\overset{d}=n-1+C_{I_n}+C'_{n-1-I_n},
 \qquad I_n\sim\operatorname{Uniform}\{0,\ldots,n-1\}.
 \label{eq:recurrence}
\end{equation}
Our companion paper proves that an ordinary lognormal earns a higher expected
log-density than an ordinary normal for every $n\geq4$. Yet, after centering
and scaling, $C_n$ approaches a non-normal Quicksort law $Q$; the fitted
lognormal approaches a normal shape. A curve can therefore win at each finite
size under one score and still miss the eventual shape.

The next question is how that shape develops. Write
$\sigma^2=7-2\pi^2/3$ and
$g(u)=1+2u\log u+2(1-u)\log(1-u)$. Starting at $\nu_0=\delta_0$,
apply the limiting split transform
\begin{equation}
 \nu_{k+1}=\mathcal T(\nu_k),\qquad
 \mathcal T(\nu)=\mathcal L\!\left(
 U X_1+(1-U)X_2+g(U)/\sigma\right),
 \quad X_1,X_2\stackrel{\mathrm{iid}}{\sim}\nu,\ U\sim{\rm Unif}(0,1).
 \label{eq:transform}
\end{equation}
These centered laws approach $Q$, and
$\operatorname{Var}(\nu_k)=r_k:=1-(2/3)^k$. A fractional generation is
$\nu_{k,\alpha}=(1-\alpha)\nu_k+\alpha\nu_{k+1}$ with variance
$r_{k,\alpha}=(1-\alpha)r_k+\alpha r_{k+1}$.
For input size $n$, our smooth count model is
\begin{equation}
 X_{n,k,\alpha}=\mu_n+
 s_n\,V_{k,\alpha}/\sqrt{r_{k,\alpha}},
 \quad V_{k,\alpha}\sim\nu_{k,\alpha},\quad
 \mu_n=\mathbb EC_n,\quad s_n^2=\operatorname{Var}(C_n).
 \label{eq:interpolation}
\end{equation}
It has the exact mean and variance of $C_n$, while its shape changes with
recursive generation. The fitted interpolation describes central probabilities
well but understates early-size skewness and the long right tail.

\section*{A finite bit measure for the count distribution}
Write $p_n(c)=\Pr(C_n=c)$ for the exact probability of integer count $c$.
Let $F_{n,\theta}$ be the CDF of a continuous approximation expressed in
\emph{comparison-count units}. Ask the same question of both models by rounding
the continuous output to the nearest integer:
\begin{equation}
 q_{n,\theta}(c)
 =F_{n,\theta}(c+\tfrac12)-F_{n,\theta}(c-\tfrac12).
 \label{eq:rounding}
\end{equation}
Now $p_n$ and $q_{n,\theta}$ live on the same count scale. With
$m=(p_n+q_{n,\theta})/2$, their Jensen--Shannon divergence in bits is
\begin{equation}
 \operatorname{JS}_2(p_n,q_{n,\theta})
 =\tfrac12 D_{\mathrm{KL},2}(p_n\Vert m)
 +\tfrac12 D_{\mathrm{KL},2}(q_{n,\theta}\Vert m).
 \label{eq:js}
\end{equation}
It is zero exactly when the rounded model gets every count probability right.
Operationally, it is how much one observed count reveals about whether it came
from Quicksort or the model, assuming either source was equally likely.

\begin{center}
\begin{tabular}{@{}rcc@{}}
\toprule
Input size $n$ & Fractional interpolation & Finite-size split correction \\
\midrule
16 & 0.00509 & 0.00186 \\
32 & 0.00235 & 0.00135 \\
64 & 0.00085 & 0.00070 \\
\bottomrule
\end{tabular}\\[-1pt]
\small Two independent Monte Carlo runs; Jensen--Shannon bits per rounded count.
\end{center}

\textbf{Reading the table.} At $n=16$, restoring finite-size structure cuts this mismatch by about 63\%.
The $n=16$ correction uses exact child laws through size 8; the $n=32,64$
corrections use moment-matched $Q$ child laws. The numbers are exploratory
Monte Carlo estimates, especially close at $n=64$.

\section*{An exact closure that teaches us what the score means}
Let $U_n\sim{\rm Unif}(-\tfrac12,\tfrac12)$ be independent of the exact count
$C_n$, and put $Y_n=C_n+U_n$. This $Y_n$ has a continuous density, yet
\begin{equation}
 \operatorname{round}(Y_n)=C_n\quad\hbox{almost surely},\qquad
 q^Y_n(c)=p_n(c),\qquad \operatorname{JS}_2(p_n,q^Y_n)=0.
 \label{eq:lift}
\end{equation}
Indeed, each interval $(c-\tfrac12,c+\tfrac12)$ is mapped back to $c$.
Thus a continuous law can match every discrete count probability at the
chosen resolution. This construction copies the exact law and explains none
of it. The useful goal is a compact model that predicts \emph{unseen} sizes.
Without a shared rounding rule, ordinary KL divergence between a discrete
law and a continuous one is infinite, regardless of how closely their CDFs
track each other.

\section*{A structural route to closing the useful gap}
Let $Z_m$ be independent, centered, unit-variance approximations to the
standardized child count laws, and set $s_m=0$ when $m\leq2$. Restore the
finite-size split $I\sim{\rm Unif}\{0,\ldots,n-1\}$:
\begin{equation}
 \widetilde C_n=n-1+\mu_I+s_I Z_I+
 \mu_{n-1-I}+s_{n-1-I}Z'_{n-1-I}.
 \label{eq:splitmodel}
\end{equation}
Conditioning on $I$ and using the exact mean and variance recurrences gives
\begin{equation}
 \mathbb E\widetilde C_n=\mu_n,\qquad
 \operatorname{Var}(\widetilde C_n)
 =\operatorname{Var}(\mu_I+\mu_{n-1-I})
   +\mathbb E(s_I^2+s_{n-1-I}^2)=s_n^2.
 \label{eq:moments}
\end{equation}
This moment guarantee does not force the tails to match. The table shows the
gain when $Z_m$ is based on $Q$, with exact small children at $n=16$.

There is also a simple error guarantee for a proposed \emph{count-level}
version. Let $P_m$ be the exact child law and $R_m$ any discrete approximation
(for example, a rounded smooth child law). Form $\widetilde P_n$ from
\eqref{eq:recurrence} with $R_i,R_{n-1-i}$ in place of the true child laws.
Convolution cannot increase total variation, so
\begin{equation}
 d_{\rm TV}(P_n,\widetilde P_n)
 \leq \frac{2}{n}\sum_{m=0}^{n-1}d_{\rm TV}(P_m,R_m).
 \label{eq:tvbound}
\end{equation}
This bound concerns the count-level variant; the table used continuous
children followed by rounding the parent. Both suggest the next test:
improve child laws by size, then evaluate the resulting parent at held-out
sizes.

\newpage

\section*{A different loss: forgetting the growing tree}
The same comparison count can arise from different binary search trees.
Under the natural growing-tree coupling, let $K_n(d)$ be the number of empty
insertion slots at depth $d$, and let $D_{n+1}$ be the depth chosen next. Then
\begin{equation}
 C_n=\sum_d dK_n(d)-2n,
 \qquad \Pr(D_{n+1}=d\mid K_n)=\frac{K_n(d)}{n+1}.
 \label{eq:profile}
\end{equation}
Knowing only $C_n$ is like knowing the total cost while forgetting where the
remaining slots are. Two trees with the same cost can give different chances
for the next insertion depth. The exact expected predictive penalty, in bits,
for retaining only cost is
\begin{equation}
 L_n=I(D_{n+1};K_n\mid C_n)
 =H(D_{n+1}\mid C_n)-H(D_{n+1}\mid K_n).
 \label{eq:predictive}
\end{equation}
Equivalently, $L_n$ is the smallest expected extra log-loss of any predictor
that sees only $C_n$, relative to one that sees $K_n$:
\begin{equation}
 \inf_{a(\cdot\mid C_n)}
 \mathbb E\,D_{\mathrm{KL},2}\!
 \left(\Pr(D_{n+1}\in\cdot\mid K_n)\,\middle\|\,a(\cdot\mid C_n)\right)
 =L_n.
 \label{eq:optimalprediction}
\end{equation}
The minimizing predictor is $\Pr(D_{n+1}\in\cdot\mid C_n)$.
This result concerns the natural tree coupling, not a particular smooth curve.

Exact enumeration gives $L_{16}=0.31545$ bits. The largest value through
$n=26$ occurs at $n=14$, about $0.31866$ bits, followed by a decline to
$L_{26}=0.27815$. Two independent large-tree simulations estimate roughly
$0.06$ bits at $n=512$. The evidence points toward a shrinking one-step gap,
but does not prove that it vanishes. Meanwhile, the number of hidden profile
details can grow: $H(K_n\mid C_n)$ rises from $2.38$ bits at $n=14$ to
$5.85$ bits at $n=26$. More of the tree is hidden by the total cost, yet less
of that hidden detail helps predict the very next insertion.

\section*{What ``close the gap'' means}
The two bit values answer different questions. Equation~\eqref{eq:js} tests
\emph{one-time distributional fit} after declaring the count resolution.
Equation~\eqref{eq:predictive} tests \emph{next-step prediction} after
compressing the tree to one number. They cannot be added or compared as parts
of a single information budget. The exact continuous lift in
\eqref{eq:lift} makes the distinction formal: because $C_n$ can be recovered
from $Y_n$ and $U_n$ supplies no tree information,
\begin{equation}
 \operatorname{JS}_2(p_n,q^Y_n)=0
 \quad\text{while}\quad
 I(D_{n+1};K_n\mid Y_n)=I(D_{n+1};K_n\mid C_n)=L_n.
 \label{eq:separation}
\end{equation}
At $n=16$, the first quantity is zero and the second is $0.31545$ bits.

For the first goal, take each $Z_m$ from the standardized fractional family
at generation $t_\beta(m)$, then let $q_{n,\beta}$ be the rounded parent law
produced by \eqref{eq:splitmodel}. Fit the shared size-to-generation map on
some sizes and test it on others:
\begin{align}
 \widehat\beta
 &=\arg\min_\beta \frac{1}{|\mathcal N_{\rm train}|}
   \sum_{n\in\mathcal N_{\rm train}}
   \operatorname{JS}_2(p_n,q_{n,\beta}), \nonumber\\
 \mathcal E_{\rm test}(\widehat\beta)
 &=\frac{1}{|\mathcal N_{\rm test}|}
   \sum_{n\in\mathcal N_{\rm test}}
   \operatorname{JS}_2(p_n,q_{n,\widehat\beta}).
 \label{eq:heldout}
\end{align}
Skewness and right-tail probabilities should be checked alongside this score.
For dynamics, the open theoretical question is the large-$n$ behavior of
$L_n$. The key modeling choice is how much \emph{state} to retain, alongside
how closely a smooth law fits the count distribution.

\vfill
{\footnotesize\textbf{Reproducibility.} The companion paper source,
\texttt{quicksort\_*.py} experiment scripts, and dated \texttt{artifacts/}
notes in this repository contain the recurrence calculations, simulations,
and limitations behind these figures.}
\end{document}
